\chapter{Data}
\label{chap:data}

\section{Executive summary}

This chapter describes the dataset used to construct and evaluate the closed-pool synonym-ranking system: the origin and documented provenance of the public source; how the project-specific dataset and its closed candidate vocabulary were built; the unit of observation and its single-label limitation; string normalization; the composition of the retained terms; and the grouped train/validation/test partitioning. Data-related limitations are discussed together with the project's other limitations in the Conclusions chapter.

The public source contains 1,266 records. After removing rows with an unusable input term, removing exact duplicate pairs, and removing one identity pair, the project retains 969 rows for evaluation. The closed candidate vocabulary of 921 unique synonym strings is built one step earlier, from the 970-row table that still includes the identity pair, so that the synonym it contributes remains a valid candidate even though the identity row itself is excluded from evaluation. These 969 rows are partitioned, by exact input term and with a fixed random seed (42), into 581 training, 194 validation, and 194 test rows.

\section{Dataset origin and scope}

The source dataset is published on Hugging Face under the identifier \texttt{\seqsplit{jensjorisdecorte/TU-Expert-Collection-Topic-Synonyms}} (Hub revision \texttt{\seqsplit{fb767656de4cb453d3b4d88bcdd060ee17a81f4b}}, last modified 2024-09-23). At this revision, the repository exposes one split, named \texttt{train}, containing 1,266 records with the following fields: an integer row index; a \texttt{topic} identifier such as \texttt{topic-12722}; an English term (\texttt{en}) and its Dutch counterpart (\texttt{nl}); and an English synonym (\texttt{en\_synonym}) with its Dutch counterpart (\texttt{nl\_synonym}).

The repository does not provide detailed documentation on the dataset's original collection or annotation procedure.

The public data is bilingual and covers heterogeneous terminology, including examples from economics, law, psychology, information systems, religion, education, and public policy. Some records have a missing, empty, or placeholder value in the English term field while the English synonym field remains populated; the exact effect of this on the retained dataset is quantified in the next section.

\section{Project-specific dataset construction}

The project retains only the two fields required by the ranking task: \texttt{en}, the input query term, and \texttt{en\_synonym}, the row-specific relevant candidate. The \texttt{topic}, \texttt{nl}, and \texttt{nl\_synonym} fields, and the numeric row index, are not passed to any evaluated method. This restricts the task to English-only synonym ranking and prevents auxiliary fields from influencing candidate generation or ranking.

The reduction from the public source to the final evaluation dataset proceeds in three filtering steps. Of the 1,266 source rows, 284 have an invalid \texttt{en} value (null, empty, or the placeholder ``\texttt{-}'') and are removed. A further 12 rows are exact duplicate (\texttt{en}, \texttt{en\_synonym}) pairs and are removed, leaving a prepared table of 970 rows. Finally, the single identity row in which the input term and its recorded synonym coincide (\emph{modernisation}--\emph{modernisation}) is removed, leaving the 969-row evaluation dataset. This identity row is removed because the retrieval implementation excludes the input term from its own candidate list; for an identity row, that exclusion would leave no correct candidate available.

Each retained row represents one evaluation instance: one input academic or domain term and one row-specific expected English synonym. The dataset is therefore not a conventional multiclass classification dataset with a small set of balanced classes; it is a ranking dataset in which each query has one labeled relevant item within a much larger shared candidate vocabulary.

For a candidate vocabulary of size $V$, each row contains one labelled relevant candidate and $V-1$ candidates treated as non-relevant under the single-label evaluation protocol used for scoring. This is a statement about the evaluation protocol, not a claim that the remaining $V-1$ candidates are necessarily semantically incorrect for that row: a candidate that is not labelled as relevant for a given row is simply not the row's recorded gold value, and may still be a plausible synonym.

This distinction matters because ten input terms occur twice in the evaluation dataset, each time paired with a different recorded synonym. Given 969 rows with ten terms occurring twice, the evaluation dataset contains 959 distinct input strings (the corresponding prepared, 970-row table contains 960 distinct input strings, one more because it still includes the removed identity row).

\section{Candidate vocabulary}
\label{sec:candidate-vocab}

The project adopts a closed-pool ranking formulation, and it is necessary to distinguish two datasets that this formulation depends on. Let the prepared dataset, obtained after removing invalid rows and exact duplicates, be

\[ D_{\mathrm{prepared}} = \{(q_i, y_i)\}_{i=1}^{970}, \]

and let the evaluation dataset, obtained after additionally removing the identity pair, be

\[ D_{\mathrm{eval}} = \{(q_i, y_i)\}_{i=1}^{969} \subset D_{\mathrm{prepared}}, \]

where $q_i$ is an English input term and $y_i$ is its recorded English synonym. The global candidate vocabulary is built from the prepared dataset, \emph{before} the identity pair is removed:

\[ C = \operatorname{unique}\left(\{y_i \mid (q_i,y_i) \in D_{\mathrm{prepared}}\}\right), \qquad |C| = 921. \]

The train, validation, and test partitions (Section~\ref{sec:partitioning}) are derived from $D_{\mathrm{eval}}$, not from $D_{\mathrm{prepared}}$. Consequently, the candidate vocabulary $C$ still contains the synonym string contributed by the excluded identity row, even though that row itself is never scored, since it does not appear in $D_{\mathrm{eval}}$ or in any split derived from it.

The vocabulary is built by deduplicating the retained synonym values under a shared normalization rule, defined in full in Section~\ref{sec:normalization}: each string is lowercased and stripped of surface punctuation before comparison, so that, for example, \emph{IT-Governance} and \emph{it governance} are treated as the same candidate. For each normalized value, the first surface form encountered is the one retained, and entries are kept in the order in which they were first observed rather than sorted alphabetically. Equality among candidates, and between a candidate and a row-specific gold value, is evaluated after this normalization, so entries differing only in case or surface punctuation are not treated as distinct. For every input term, the system produces an ordered list of candidates selected from $C$; the raw output of any language model used for query expansion is not treated as a final answer unless it corresponds to an entry of $C$, so every method, including the ones that free-generate a candidate proposal, ultimately returns a member of the closed vocabulary rather than unconstrained text.

\section{Preprocessing and normalization}
\label{sec:normalization}

\subsection{Stored representation}

The stored dataset retains the original surface form of each string after basic cleaning: leading and trailing whitespace is removed, and empty strings or the placeholder ``\texttt{-}'' are treated as missing values and dropped.

\subsection{Equality representation}

Independently of the stored representation, every equality check performed by the project -- candidate deduplication, exclusion of the input term from its own candidate list, and matching a prediction against the row-specific gold label -- uses a shared normalized form of each string: lowercased, stripped of leading and trailing whitespace, filtered to remove characters outside basic ASCII letters, digits, whitespace, and hyphens, with the remaining whitespace collapsed. This filter deletes unsupported characters rather than transliterating them to an ASCII equivalent, so an accented character is removed rather than folded to its unaccented counterpart; given that the retained fields are English-language academic terminology, this is expected to affect a small share of entries. The stored strings themselves are not modified by this process -- it is used only for equality comparisons.

\section{Dataset composition}

\subsection{Term-type distribution}

The 969-row evaluation dataset is divided into four mutually exhaustive, heuristically defined term types: abbreviation-like terms, compounds, multi-word phrases, and single words. A term is classified as abbreviation-like if it is a single token that is either fully uppercase or at most five characters long; a single token that does not meet this criterion is a single word; among multi-token terms, a term containing a hyphen or consisting of exactly two words is a compound; every remaining multi-token term is a multi-word phrase. These are project-specific analytical categories rather than authoritative linguistic classes, but the classification rule is deterministic, with no ambiguous cases. Table~\ref{tab:term-types} reports the resulting distribution.

\begingroup
\footnotesize
\begin{table}[htbp]
\centering
\begin{tabular}{lrr}
\toprule
\textbf{Term type} & \textbf{Rows} & \textbf{Dataset share}\\
\midrule
Abbreviation-like & 34 & 3.5\%\\
Compound & 432 & 44.6\%\\
Multi-word phrase & 250 & 25.8\%\\
Single word & 253 & 26.1\%\\
\midrule
Total & 969 & 100.0\%\\
\bottomrule
\end{tabular}
\caption{Descriptive term-type distribution of the 969-row evaluation dataset. Model performance broken down by these categories is reported in the Results chapter.}
\label{tab:term-types}
\end{table}
\endgroup

Compounds form the largest category, at approximately 44.6\% of the dataset; multi-word phrases and single words each represent approximately one quarter; abbreviation-like inputs form a comparatively small subset. A lexical-overlap indicator, based on shared whitespace-separated tokens between an input term and its gold synonym, was also defined for post-hoc error analysis; the corresponding performance results are reported in the Results chapter rather than here.

\section{Train, validation, and test partitioning}
\label{sec:partitioning}

The 969-row evaluation dataset was partitioned with a fixed random seed (42) and a nominal 60/20/20 allocation, producing 581 training rows, 194 validation rows, and 194 test rows. The partition is grouped by the exact input term rather than by row, so that every occurrence of a repeated input term is assigned to the same partition. This grouping is necessary because ten input terms occur twice with different recorded synonyms: without it, one occurrence of a term could be used for configuration decisions while another occurrence of the same term, with a different gold label, appeared in the test partition.

Table~\ref{tab:partitions} summarizes each partition's size and its actual role in the project.

\begingroup
\footnotesize
\begin{table}[htbp]
\centering
\begin{tabular}{p{2.2cm}p{1.2cm}p{1.3cm}p{8.6cm}}
\toprule
\textbf{Partition} & \textbf{Rows} & \textbf{Share} & \textbf{Role in this project}\\
\midrule
Train & 581 & 59.96\% & Not used to fit or select any of the reported systems.\\
Validation & 194 & 20.02\% & Used throughout the project for model, fusion, and configuration decisions.\\
Test & 194 & 20.02\% & Not used for direct parameter or model selection; several already-selected configurations were evaluated on it at different points as the experimental process progressed.\\
Full & 969 & 100.00\% & Not an independent split; the union of the three partitions, used only for descriptive reporting and error analysis after a final configuration had been selected.\\
\bottomrule
\end{tabular}
\caption{Role of each partition, as actually used in this project.}
\label{tab:partitions}
\end{table}
\endgroup

The training partition is not used to fit or select any of the evaluated systems: no component in this project is trained with gradient descent, so there is no model-fitting step for it to support. It is reserved, unused, so that a future extension of this project -- for example, fitting fusion weights or a confidence threshold on data the reported comparisons never touched -- would still have a partition available that has not influenced any decision made in this thesis. The validation partition, by contrast, was used throughout the project to choose between embedding models, fusion strategies, query-expansion settings, candidate-pool widths, and rerankers.

The test partition was not used to choose between candidate methods or parameter values -- those decisions were made on the validation partition. It was, however, evaluated more than once: as the pipeline evolved, several configurations already selected on validation were subsequently run on the test partition at different points during the project, and their results are reported alongside the final configuration in the Results chapter. The test partition was therefore not strictly held out until a single final evaluation. Even though no configuration was chosen because of its test-set score, repeated access to test-set results over the course of development can introduce indirect influence on later decisions, and this risk should be weighed when interpreting the reported test numbers.
